\summarychapter{25}{Radical Discipleship}{The Anabaptists}

\begin{inbrief}
The Anabaptists were the radicals of the Reformation, calling the moderate reformers to cut deeper.
They wanted churches of committed believers, baptized on their own confession and free of the state,
and for that they were hunted by Catholics and Protestants alike. The chapter tells of their birth in
Zurich, their persecution, the disaster at M\"unster, and the four convictions by which they became
forerunners of nearly all modern Protestants.
\question{How could people who only wanted to recover the church of the New Testament end up so far in advance of their age?}
\end{inbrief}

On the snowy night of January 21, 1525, a dozen men made their way through Zurich to the house of
Felix Manz near the Great Minster. That day the city council had ordered Conrad Grebel and Manz to
stop their Bible classes; four days earlier it had ordered every infant baptized by the eighth day, on pain of the parents' banishment. After prayer the former priest George Blaurock asked Grebel to
baptize him in the apostolic way, on confession of personal faith. Grebel did, and Blaurock baptized
the others. Anabaptism was born. Its direct descendants are the Mennonites and Hutterites, of whom
only the Old Order Amish keep to horse and buggy. What unites them is a set of beliefs, many now
shared by Baptists, Quakers and Congregationalists. Like the Benedictines, Shelley observes, they show
how often the people most devoted to the next world are the ones who alter this one.

\ssection{Basic Beliefs of Anabaptists}
Most revolutions produce radicals who want to reform the reformation. Like other countercultural
movements, the Anabaptists had no single body of doctrine and no common organization. Even the name,
``rebaptizers,'' was fastened on them by enemies in order to link them with ancient heretics and
expose them to persecution, which it did. They themselves denied rebaptizing, since they did not
regard the sprinkling of infants as baptism at all, and would have preferred to be called Baptists.
What was really at stake was what the church is and how it stands toward the civil power.

They reached their convictions through Scripture, taking Luther at his word that common people may
search the Bible for themselves. In the New Testament they found no alliance of church and state and
no Christendom, but companies of people who had each decided freely to follow Jesus---a
revolutionary idea in the sixteenth century. For all Luther's stress on personal faith, Lutheran
churches were established churches, counting the whole population as members and living on state
support; official Protestantism seemed little different from official Catholicism. The Anabaptists
aimed at the ``restitution'' of apostolic Christianity. Only those who had experienced spiritual
regeneration are fit for baptism; infant baptism was a device for perpetuating a nominal Christian
society. The true church is a community of saints, disciples in a wicked world, shaping society by
example and if need be by death. They refused to bear arms, hold political office or take oaths,
which looked to their age like sedition or treason. By mid-century there were three groups in
German-speaking Europe: the Swiss Brethren at Zurich, the Hutterites in Moravia and the Mennonites in
the Netherlands and North Germany.

\ssection{Radicals in the Swiss Alps}
Grebel and Manz, educated men of standing, began as supporters of Ulrich Zwingli (1484--1531). In
1519 Zwingli became people's priest at the Great Minster and started the reformation in Zurich not
with theses but with biblical sermons. Influenced by Erasmus, he revered the New Testament; a young
hearer, Thomas Platter, felt as if he were being lifted by the hair of his head. In one respect
Zwingli was stricter than Luther. Luther allowed whatever the Bible does not forbid; Zwingli rejected
whatever it does not prescribe. So Zurich stripped away candles, statues, music and pictures, the
spirit later called Puritanism in England.

Following Zwingli's own method, Grebel and Manz saw the difference between the apostolic churches and
their own. In Zurich, as everywhere, every child was baptized and church and society were identical,
whereas the New Testament church was a fellowship of the few, true believers ready to live and die for
their Lord. They wanted such a church, free of the state. Zwingli, who needed the city fathers,
refused. In the fall of 1524 the Grebels declined to have their newborn son baptized, and other
parents followed. After a public debate on January 17, 1525, the council declared Zwingli the winner
and gave the parents a week. The baptisms of January 21 were the answer. The little company withdrew
to the village of Zollikon, where the first Anabaptist congregation was formed, the first free church
of modern times.

The authorities arrested and imprisoned them, but on release they sought converts in neighboring
towns. On March 7, 1526, the council decreed death by drowning for anyone who rebaptized: heretics who wanted water were to have it. On January 5, 1527, Felix Manz was drowned in the Limmat, the
first Anabaptist martyr, and within four years the movement around Zurich was nearly wiped out. Those
who fled to Germany and Austria fared no better. In 1529 the imperial Diet of Speyer declared
Anabaptism a heresy punishable by death, and during the Reformation years four to five thousand died by fire, water or the sword. To us they seem to have asked only for the right to their own
beliefs; to their age they seemed to destroy the fabric of society. The chapter quotes a letter of
1573 from a young mother in an Antwerp jail, spared just until her child was born, to her daughter born days before. Her husband had been put to death earlier. She begs the child not to be ashamed
of her parents, for theirs is the way the prophets and apostles went.

George Blaurock carried the message east to the Tyrol, where he was burned on September 6, 1529.
Persecution drove many north to Moravia, where tolerant princes sheltered them. There they founded
the \emph{Bruderhof}, a Christian commune modeled on the apostolic community and useful for survival,
to show that in God's kingdom brotherhood comes before self. Under Jakob Hutter, who died in 1536,
these groups became the Hutterites.

\ssection{Kingdom Building Gone Mad}
Fear of the radicals deepened with the rebellion at M\"unster, a bishop's city in Westphalia. The
Reformation spread there in 1532, at first in a conservative Lutheran form. Then immigrants, apostles
of Jan Matthijs, brought fanaticism, and many expected the Lord's earthly kingdom to be set up in the
city; such belief in a thousand-year earthly reign of Christ is called chiliasm. When the bishop
besieged the city, these Anabaptists, uncharacteristically, took up arms. In mid-1534 Jan of Leiden, who had kept an inn, took control, claimed new revelations, introduced polygamy and took
the title ``King David.'' He lived in splendor with his harem and kept up morale in spite of hunger
until the city fell on June 24, 1535. For centuries afterward the word Anabaptist meant M\"unster and
wild-eyed fanaticism to Europeans.

The dispirited Anabaptists of the lower Rhine took heart again through Menno Simons (about
1496--1561), a former priest who, always in danger, traveled among the scattered groups of northern
Europe and preached to them by night. He was unbending in requiring pacifism, and though he did not
found the movement, most descendants of the Anabaptists bear his name: Mennonites.

\ssection{Pioneers of Modern Christianity}
Living as outlaws, the Anabaptists had little chance to state their beliefs together. They did so
once. In 1527 they met at Schleitheim, on today's Swiss-German border, in the first synod of the
Protestant Reformation, led by the former Benedictine Michael Sattler, who was burned four months
later. Their ``Brotherly Union,'' the Schleitheim Confession, came to be accepted by most Anabaptists
within a decade. John H.~Yoder and Alan Kreider summarize it in four convictions.

The first is discipleship. The Christian's relation to Christ must go beyond inner experience and
doctrine to a daily walk in which his teaching and example transform life: ``No one can truly know
Christ except he follow him in life.'' Where most Christians think that one must understand before
obeying, the Anabaptist instinct is that only obedience brings understanding. Hence, for example,
their refusal of oaths, because Jesus said, ``Do not swear at all''; truth has no degrees. The second
is love. Toward outsiders they were pacifists who would not go to war, defend themselves against
persecutors or share in the state's coercion; among themselves they practiced mutual aid and the
redistribution of wealth, in Moravia even communal living. The third is the congregational view of
authority, toward which Luther and Zwingli had leaned in their first years. All members are believers
baptized voluntarily; each is a priest to fellow believers and a missionary to unbelievers. Decisions
rest with the whole membership, Scripture is interpreted by the consensus of the local gathering,
and discipline is exercised together. The fourth is the separation of church and state. Christians
are a ``free, unforced, uncompelled people.'' Faith is God's gift, and rulers exceed their competence
when they ``champion the Word of God with a fist.'' The church is distinct from society even when
society calls itself Christian. It is a pilgrim people. The Anabaptists were thus the first in modern
times to preach thorough religious liberty, meaning the freedom to worship with fellow believers while the state neither funds nor punishes it. Their descendants lost much of this over the centuries. The search for a pure church often turned
legalistic, and for the sake of survival they gave up evangelism and became known as good farmers,
the ``Quiet in the Land.'' Revival and growth came only in the late nineteenth century, and by the
late twentieth their worldwide membership passed half a million. Far beyond their own communities
Christians now hold one or more of the principles for which the first radicals were willing to die.

\begin{keyfacts}
\keypeople{Conrad Grebel; Felix Manz; George Blaurock; Ulrich Zwingli; Erasmus; Jakob Hutter; Jan
Matthijs; Jan of Leiden; Menno Simons; Michael Sattler.}
\keydates{%
  \kd{1519}{Zwingli begins preaching reform at the Great Minster in Zurich}
  \kd{1525}{First believers' baptisms in Zurich (January 21); first free church at Zollikon}
  \kd{1527}{Felix Manz drowned (January 5); Schleitheim Confession}
  \kd{1529}{Diet of Speyer makes Anabaptism punishable by death; Blaurock burned}
  \kd{1534--1535}{The M\"unster rebellion under Jan of Leiden}
  \kd{1536}{Death of Jakob Hutter, organizer of the Hutterites}
  \kd{1561}{Death of Menno Simons}}
\keyterms{Anabaptist (``rebaptizer''); believers' baptism; restitution; free church; Swiss Brethren;
Hutterites; \emph{Bruderhof}; Mennonites; chiliasm; discipleship; pacifism; congregational authority;
separation of church and state; religious liberty.}
\end{keyfacts}
